\section{Conclusiones}
    Realizar este proyecto me permitió profundizar acerca de los beneficios prácticos 
    de la FFT, como la reducción del costo computacional y filtrado de ruido senoidal, 
    pero también sobre algunas de las limitaciones que tienen los filtros de recorte 
    basados en el de Butterworth en ciertos escenarios. Lo relevante de este proyecto 
    se encuentra en que si queremos aplicar un filtro de convolución a una imagen 
    muy grande, lo más conveniente sería aplicarlo en el espacio de frecuencias debido 
    a que se resume en una multiplicación matricial, este cambio operacional reduce 
    drásticamente el costo computacional cuando se aplica el mismo filtro sobre miles 
    de imágenes, esto en el caso del entrenamiento de una CNN.\\

    Existen muchos tipos diferentes de ruido con una diversidad de fuentes, algunos de 
    estos ruidos se pueden mitigar con las técnicas vistas previamente y más si 
    son localizados, puntuales o posibles de suavizar. Pero qué ocurre cuando el ruido 
    está presente de forma uniforme sobre toda la imagen, como estática o en senoidales, 
    para eliminarlo requerimos usar de otras técnicas que permitan observarlo como 
    si fuera un punto aislado o diferente, de aquí que usar la FFT nos genere esta 
    representación que estamos buscando, y esto se apoya de los resultados obtenidos 
    de forma experimental donde cada destello raro tiene el potencial de ser parte 
    del ruido senoidal presente en las imágenes.\\ 

    Por lo anterior, se tiene que probar y experimentar con diferentes filtros para lograr 
    el efecto deseado en la imagen resultante, en algunos casos usar los de recorte 
    de Butterworth no son los más convenientes debido a su efecto puntual. Como se mostró 
    en algunas de las imágenes, el ruido también puede aparecer como líneas rectas entonces 
    ahí se requiere de otros tipos de filtros que actúen sobre estos elementos o geometrías. 